\scanpage{97}
6) \uline{\unsure{formes|forme} quadratiques \unsure{invariante|invariantes}}.

On s'intéresse aux expressions
\[
\deleted{f}F(x,y)=ax^2+bxy+cy^2+ux+vy+d
\]
polynômiales de degré $\le2$ sur $E$, qui correspondent aux formes quadratiques sur l'enveloppe vectorielle $\unsure{\mathcal E|\mathscr E|\hat E}$ de $E$. On cherche celles qui sont invariantes par $G=\mathrm{Aut}\,\Pi$, i.e. par $\sigma_0,\sigma_1$. Il y a les constantes, donc on est ramené à voir les conditions pour que \deleted{$f$ soit \unsure{inv.}} \inserted{l'expression sans} terme constant
\[
\deleted{f}F(x,y)=ax^2+bxy+cy^2+ux+vy
\]
soit invariante. L'invariance par $\sigma_1$ signifie que $f(x,y)$ est symétrique en $x,y$
\[
\begin{aligned}
a=c,\quad u=v\qquad \text{i.e.}\\
\deleted{f}F(x,y)=a(x^2+y^2)+bxy+u(x-y)
\end{aligned}
\]
et quand on explicite l'invariance par $\sigma_0$, on trouve
\[
\deleted{f}F=a f_{\unsure{\alpha|*}}
\]
\[
\boxed{f_\alpha(x,y)=(x^2-x)+(y^2-y)-\alpha xy}
\]
Les fonctions $f$ \inserted{(pol. de degré $\le\unsure{2|?}$)} nulles sur $S$ invariantes par $G$ sont les multiples \unsure{de|du} $f_\alpha$ ci-dessus.
% margin note written sideways
\margin{NB}
% archive sheet number: 96
